\subsection{Reproduktion und Status}\label{reproduktion-und-status}

\codeword{contracted\_krylov.cpp} ist der unabhängige ganzzahlige
Kontraktionskern. \codeword{verify.py} prüft diesen Kern, die
korrigierte Pythagoras-Rechnung, kleine Ritz-Matrizen, sämtliche
angeführten Graph-, Kern- und Gradierungsidentitäten sowie die
Matrixalgebra-Zertifikate. Die Ergebnisse stehen in
\codeword{audit.json} und \codeword{audit\_optimized.json}.

Es bestehen \textbf{376 Prüfbedingungen: 371 exakte und fünf
numerische}. Für eine eigenständige Reproduktion zuerst den kleinen
C++-Kern mit C++17 und Optimierung übersetzen; anschließend
\codeword{verify.py} normal und mit \codeword{-OO} starten. Sämtliche
Eingaben liegen eingefroren unter \codeword{sources/}; die
Ergebnisdateien enthalten ihre SHA-256-Werte. Fremde Pythonprogramme
werden hierbei nicht ausgeführt. Die numerischen Ritz-Prüfungen sind als
solche markiert, alle Algebra-Zertifikate und die vollständige direkte
Kontraktion verwenden exakte Ganzzahl-/Bruchrechnung.

Beide abschließenden Ausführungen bestanden und lieferten byteidentische
JSON-Dateien mit SHA-256
\codeword{148047028236a887de1b28a0a9c12d84037a7c633892abc51dd9047f0b7a7baf}.

Keine T1-T8-Abschlussbehauptung. Keine Änderungen an fremden Quellen,
Hauptpaper, Webseite, Ledger oder Git-Historie.

\begin{center}\rule{0.5\linewidth}{0.5pt}\end{center}

\section{Anhang C: minimaler Transfer und seine
Quellen-Grenze}\label{anhang-c-minimaler-transfer-und-seine-quellen-grenze}

\section{Minimaler Ursprungsstrang: geteilte Moden und ein echter
nativer
Vierzustandskanal}\label{minimaler-ursprungsstrang-geteilte-moden-und-ein-echter-nativer-vierzustandskanal}

Stand: 15. September 2026. Unabhängiger, begrenzter Forschungsstrang.

\subsection{Ergebnis und Reichweite}\label{ergebnis-und-reichweite}

Es gibt zwei verschiedene Ergebnisse, die ausdrücklich nicht verwechselt
werden dürfen:

\begin{enumerate}
\def\labelenumi{\arabic{enumi}.}
\tightlist
\item
  \textbf{Ein exakt lösbarer Überlappungsbaustein:} Eine gemeinsame
  Fermionmode und ein gemeinsamer Bosonkanal werden in jedem festen
  positiven Ressourcensektor zu einer gewöhnlichen effektiven
  Fermionmode. Dafür ist keine direkte Endpunkt-Hopping-Wechselwirkung
  nötig. Dieser Dreierstern ist jedoch \textbf{kein Ausschnitt einer
  tatsächlichen nativen W-Zeile}.
\item
  \textbf{Ein tatsächlicher W-basierter Vierzustandskanal:} Im
  unveränderten 64-Fermion/60-Boson-Paartensor existiert ein exakt
  invariantes Vierzustandsystem, wenn ein \textbf{zusätzlicher reiner
  Bosonmischer} zwischen Kanal 0 und 1 zugelassen und ein bestimmter
  Zustand mit Gesamtladung 9 präpariert wird. Der vollständige native
  Hamiltonoperator plus dieser Mischer überträgt dort eine innere
  Fermionmarke mit einer streng abgesicherten Wahrscheinlichkeit von
  \textbf{mehr als 99,3 \%}.
\end{enumerate}

Das zweite Ergebnis benötigt keine Projektion, die die übrigen nativen
Zustände künstlich wegschneidet: Der Vierzustandsraum ist unter allen 60
ursprünglichen Paaroperatoren plus dem Mischer exakt invariant. Es ist
aber \textbf{weder eine Untersuchung der Entnahmeantwort des
N=64-Grundzustands noch ein Transport zwischen zwei unabhängigen
Banken}. Die Herkunft des Bosonmischers, die spezielle Präparation und
die räumliche Bedeutung der Marken bleiben offen.

\subsection{1. Warum eine Clockphase die bisherige Paritätsschranke
nicht
aufhebt}\label{warum-eine-clockphase-die-bisherige-parituxe4tsschranke-nicht-aufhebt}

Seien zwei operational unabhängig definierte Banken mit lokalen
Paritäten

\[\adjustbox{max width=\linewidth}{$\displaystyle 
\Pi_x=(-1)^{N_{f,x}},\qquad \Pi_y=(-1)^{N_{f,y}}
$}\]

gegeben. Kommutiert jede ausführbare lokale Operation und jeder einzelne
realisierte Messzweig mit beiden Paritäten, gilt dies auch für beliebige
Produkte, Summen, Adjungierte und starke beschränkte Grenzwerte. Eine
Clockoperation, die lediglich innerhalb jeder Bank Fermionmoden
zahlenerhaltend permutiert, bleibt in dieser Algebra.

Ein echter Einfermiontransfer erfüllt dagegen

\[\adjustbox{max width=\linewidth}{$\displaystyle 
\Pi_x f_y^\dagger f_x=-f_y^\dagger f_x\Pi_x,
\qquad
\Pi_y f_y^\dagger f_x=-f_y^\dagger f_x\Pi_y.
$}\]

Er kann daher nicht allein durch mehr Kompositionen jener Operationen
entstehen. Eine bloß paritätskovariante CP-Abbildung ist nicht die
gleiche Voraussetzung: Sie kann ungerade Krauszweige besitzen. Die
Aussage gilt nur für den ausdrücklich geraden ausführbaren
Operationssatz.

Bei geteilten Fermionmoden sind die beiden lokalen CAR-Algebren von
Anfang an nicht unabhängig. Liegt dieselbe Mode s in beiden, kann sie
nicht zugleich als zwei unabhängige antikommutierende Kopien behandelt
werden: \(\{s,s^\dagger\}=1\). Auch lokale Ladungen, die s doppelt
mitzählen, sind keine unabhängigen additiven Bankladungen. Überlappung
widerlegt somit den Paritätssatz nicht, sondern ändert seine
Voraussetzungen.

\subsection{2. Exakter CAR-Baustein aus einem geteilten Boson und
Fermion}\label{exakter-car-baustein-aus-einem-geteilten-boson-und-fermion}

Betrachtet werden drei CAR-Moden \(f_x,s,f_y\), ein CCR-Boson b und

\[\adjustbox{max width=\linewidth}{$\displaystyle 
Q_x=b^\dagger s f_x,\qquad Q_y=b^\dagger s f_y,
\qquad
H_\star=\Delta N_b+g(Q_x+Q_y+Q_x^\dagger+Q_y^\dagger).
$}\]

Der Ressourcenzähler

\[\adjustbox{max width=\linewidth}{$\displaystyle 
K=N_b+n_s
$}\]

kommutiert mit allen diesen Operationen. Auf jedem exakten K-Sektor mit
ganzzahligem \(K\ge1\) setze

\[\adjustbox{max width=\linewidth}{$\displaystyle 
c^\dagger=\frac{b^\dagger s}{\sqrt K},\qquad
c=\frac{s^\dagger b}{\sqrt K}.
$}\]

Direkt aus CCR und CAR folgen

\[\adjustbox{max width=\linewidth}{$\displaystyle 
(c^\dagger)^2=c^2=0,\qquad
\{c,c^\dagger\}=\frac{N_b+n_s}{K}=1.
$}\]

Die Mode c antikommutiert mit beiden Endpunktmoden. Ferner

\[\adjustbox{max width=\linewidth}{$\displaystyle 
N_b=K-1+n_c,
\qquad
N_f+2N_b=2K-1+n_x+n_c+n_y.
$}\]

Damit gilt exakt

\[\adjustbox{max width=\linewidth}{$\displaystyle 
\boxed{
H_\star=\Delta(K-1)+\Delta n_c
 +g\sqrt K\big(c^\dagger f_x+c^\dagger f_y+\mathrm{h.c.}\big).
}
$}\]

Die gewöhnliche Drei-Moden-Transferkette wurde hier aus einer
überlappenden Paarumwandlung erhalten. Ein nützlicher Operatorausdruck
derselben Tatsache ist

\[\adjustbox{max width=\linewidth}{$\displaystyle 
\boxed{[Q_y^\dagger,Q_x]=K f_y^\dagger f_x.}
$}\]

Das ist zunächst eine Operatoridentität. Die Verfügbarkeit getrennter
Kommutator-Kontrollsequenzen folgt daraus nicht automatisch. Schon der
konstante Summen-Hamiltonoperator zeigt jedoch Transfer. Der minimale
positive Ressourcensektor K=1 benötigt am Anfang eine besetzte
gemeinsame Fermionmode s und ein leeres Boson.

\subsubsection{Vollständige Transferrechnung im
Vergleichsmodell}\label{vollstuxe4ndige-transferrechnung-im-vergleichsmodell}

Im effektiven Einteilchenraum und bei K=1 lautet die Matrix

\[\adjustbox{max width=\linewidth}{$\displaystyle 
\begin{pmatrix}0&g&0\\g&\Delta&g\\0&g&0\end{pmatrix}.
$}\]

Der antisymmetrische Endpunktzustand ist dunkel. Der symmetrische
koppelt mit \(\sqrt2g\) an c.~Für \(g/\Delta=1/20\), \(m=101\) und

\[\adjustbox{max width=\linewidth}{$\displaystyle 
t=\frac{202\pi}{\Delta\sqrt{51/50}}
$}\]

verschwindet die mittlere Besetzung wieder exakt. Die
Zielwahrscheinlichkeit ist

\[\adjustbox{max width=\linewidth}{$\displaystyle 
p_\star=\sin^2\left[\frac\pi2\,101\left(1-\sqrt{50/51}\right)\right].
$}\]

Durch Quadrieren rationaler Schranken erhält man

\[\adjustbox{max width=\linewidth}{$\displaystyle 
\frac{199}{200}<101\left(1-\sqrt{50/51}\right)<1.
$}\]

Mit \(|\sin u|\le|u|\) und \(\pi<355/113\) folgt

\[\adjustbox{max width=\linewidth}{$\displaystyle 
p_\star>1-\left(\frac{355/113}{400}\right)^2>0.9999.
$}\]

Diese Zahl gehört \textbf{nur zum Vergleichsmodell mit gemeinsamem b}.

\subsubsection{Die direkte
Quellen-Grenze}\label{die-direkte-quellen-grenze}

Die tatsächliche native W-Zeile enthält jeweils acht disjunkte
Fermionpaare, also 16 verschiedene Marken. Eine einzelne Zeile enthält
deshalb niemals zugleich \(s f_x\) und \(s f_y\) mit \(x\ne y\). Das ist
an allen 60 Zeilen exakt geprüft.

Ein gemeinsamer-b-Dreierstern darf somit nicht allein wegen der Form
\(b^\dagger f f\) als nativer Baustein bezeichnet werden. Insbesondere
wäre ein Wechsel von zwei Bosonkanälen zu einer gemeinsamen Mode ohne
Kontrolle des orthogonalen Kanals eine zusätzliche Modellannahme.

\subsection{3. Exakter Vierzustandskanal mit dem vollständigen nativen
W}\label{exakter-vierzustandskanal-mit-dem-vollstuxe4ndigen-nativen-w}

Die gepinnte Quelle ist

\codeword{universalraum-native-ground-response-20260915/ground\_replay/outputs/simple\_core/spinor\_tensors.npz},

SHA-256
\codeword{3f00a0892157a691e4ef26920c3702772c7bf76a250fad98ff04a079bb64b763}.

Die zwei relevanten Quellenzeilen sind

\[\adjustbox{max width=\linewidth}{$\displaystyle 
\begin{aligned}
P_0={}&-f_{57}f_4+f_{56}f_5+f_{53}f_8-f_{52}f_9
       -f_{45}f_{16}+f_{44}f_{17}+f_{33}f_{28}-f_{32}f_{29},\\
P_1={}&-f_{58}f_4+f_{56}f_6+f_{54}f_8-f_{52}f_{10}
       -f_{46}f_{16}+f_{44}f_{18}+f_{34}f_{28}-f_{32}f_{30}.
\end{aligned}
$}\]

Die Konvention ist \(P_A=\sum_{i<j}W_{A,ij}f_jf_i\). Das gemeinsame
aktive Fermion ist s=4; die Endpunktmarken sind 57 und 58.

Präpariert werden sieben unveränderte Pauli-Blocker

\[\adjustbox{max width=\linewidth}{$\displaystyle 
S=\{8,16,28,32,44,52,56\}.
$}\]

Sei \(|S\rangle\) der geordnet erzeugte Fockzustand dieser Moden. Die
vier Basiszustände sind

\[\adjustbox{max width=\linewidth}{$\displaystyle 
\begin{array}{c|c|c}
\text{Zustand}&\text{besetzte Fermionmoden}&\text{Bosonen}\\\hline
L&S\cup\{4,57\}&0\\
B_0&S&1\text{ in Kanal }0\\
B_1&S&1\text{ in Kanal }1\\
R&S\cup\{4,58\}&0
\end{array}
$}\]

Jeder Zustand besitzt die native Gesamtladung \(N=N_f+2N_b=9\).

Zugelassen wird nun genau eine zusätzliche Hamiltonoperation,

\[\adjustbox{max width=\linewidth}{$\displaystyle 
H_J=J(b_1^\dagger b_0+b_0^\dagger b_1).
$}\]

Alle ursprünglichen 60 nativen Paaroperatoren bleiben erhalten:

\[\adjustbox{max width=\linewidth}{$\displaystyle 
H=\Delta N_b+g\sum_{A=0}^{59}(b_A^\dagger P_A+P_A^\dagger b_A)+H_J.
$}\]

Der Prüfer wendet sämtliche 480 signierten Paarkanäle einschließlich
aller CAR-Vorzeichen und der Bosonfaktoren auf jeden der vier Zustände
an. Es entstehen \textbf{keine Zustände außerhalb ihres linearen
Spanns}. Die Pauli-Blocker verhindern insbesondere die sieben
unerwünschten Rückpaarungen jedes aktiven Bosonkanals.

Es ergibt sich exakt

\[\adjustbox{max width=\linewidth}{$\displaystyle 
\boxed{
H_4=\begin{pmatrix}
0&g&0&0\\
g&\Delta&J&0\\
0&J&\Delta&g\\
0&0&g&0
\end{pmatrix}_{(L,B_0,B_1,R)}.
}
$}\]

Ohne Mischer, J=0, zerfällt die Matrix in zwei getrennte 2×2-Blöcke.
Dann ist \(\langle R|e^{-itH}|L\rangle\) zu jeder Zeit exakt null. Mit J
ungleich null beginnt die Transferamplitude bei dritter Ordnung; die
Wahrscheinlichkeit hat den führenden Term \(g^4J^2t^6/36\).

\textbf{Die Quelle enthält hier also die beiden Endstücke und eine exakt
funktionierende Pauli-Blockierung, aber nicht die geprüfte
Bosonverbindung als bereits verfügbare Operation.}

\subsection{4. Strenge vollständige Transfergrenze im tatsächlichen
W-Zeugen}\label{strenge-vollstuxe4ndige-transfergrenze-im-tatsuxe4chlichen-w-zeugen}

Unter Spiegelung L↔R und B₀↔B₁ zerfällt H₄ in

\[\adjustbox{max width=\linewidth}{$\displaystyle 
H_+=\begin{pmatrix}0&g\\g&\Delta+J\end{pmatrix},\qquad
H_-=\begin{pmatrix}0&g\\g&\Delta-J\end{pmatrix}.
$}\]

Wähle den vorhandenen Prüfquotienten \(g/\Delta=1/20\) und die
ausdrücklich neu gesetzte Mischerstärke

\[\adjustbox{max width=\linewidth}{$\displaystyle 
J=\Delta-\frac{2g}{\sqrt3}
 =\Delta\left(1-\frac1{10\sqrt3}\right),
\qquad t=\frac{\pi\sqrt3}{g}=\frac{20\pi\sqrt3}{\Delta}.
$}\]

Da \(0<J<\Delta\), bleiben die Eigenfrequenzen des Bosonmischers
\(\Delta\pm J\) positiv. Es wurde \textbf{keine Nullfrequenz durch
exakte Aufhebung von \(\Delta N_b\)} erzeugt. Auch auf dem ganzen
Fockraum bleibt die hinzugefügte Hamiltonfamilie nach unten
beschränkbar: Das Bosonquadrat ist strikt positiv und die endlichen
Fermionoperatoren können durch quadratische Ergänzung kontrolliert
werden. Das bedeutet nicht, dass ihr Grundzustand noch der unveränderte
native Grundzustand ist.

Für H₋ ist \(\Delta-J=2g/\sqrt3\). Zu der gewählten Zeit kehrt dessen
Endpunktkomponente exakt mit Phase −1 zurück, ohne mittlere
Restbesetzung.

Für H₊ setze

\[\adjustbox{max width=\linewidth}{$\displaystyle 
a=\Delta-\frac g{\sqrt3},\quad
\omega=\sqrt{a^2+g^2},\quad
w_- =\frac{1-a/\omega}{2}.
$}\]

Die Endpunktamplitude lautet

\[\adjustbox{max width=\linewidth}{$\displaystyle 
A_+=(1-w_-)e^{i(\omega-a)t}+w_-e^{-i(\omega+a)t}.
$}\]

Mit \(\omega-a\le g^2/(2a)\), \(w_-\le g^2/(4a^2)\),
\(\cos u\ge1-u^2/2\) und \(|A_+|^2\ge1-g^2/a^2\) folgt

\[\adjustbox{max width=\linewidth}{$\displaystyle 
\begin{aligned}
p_{L\to R}
 &=\left|\frac{A_++1}{2}\right|^2\\
 &\ge1-\frac{g^4t^2}{16a^2}-\frac{g^2}{2a^2}\\
 &=1-\frac{3\pi^2+8}{6400(a/\Delta)^2}.
\end{aligned}
$}\]

Aus \(\sqrt3>17/10\) folgt \(a/\Delta>33/34\). Somit ergibt sich rein
rational

\[\adjustbox{max width=\linewidth}{$\displaystyle 
\boxed{
p_{L\to R}>
1-\frac{3(355/113)^2+8}{6400(33/34)^2}
=\frac{2009992727}{2022609600}
>0.9937620819>0.993.
}
$}\]

Die Aussage ist eine \textbf{analytische Schranke der vollständigen
Dynamik auf einem exakt invarianten Quellen-Unterraum}. Es gibt hier
keinen numerisch abgeschnittenen Restzustandsraum. Die
64-Fermion/60-Boson-Quelle ist nicht durch vier frei erfundene
Matrixeinträge ersetzt worden: Ihre exakte Einschränkung wurde zuerst
vollständig geprüft.

Die Zeit in Einheiten \(\hbar=1\) ist ungefähr \(108.83/\Delta\). Diese
numerische Orientierung und die gesetzte Mischerstärke sind keine
vorhergesagten Naturkonstanten.

